\section{Appendix A: figure compendium}
\label{sec:appendix-figures}

This appendix reproduces the remaining figures from the Mermaid catalog
(\texttt{adoption/mermaid/}) as colored TikZ, so the framework carries the full set
of exhibits in the strict palette. 

\begin{cfwfig}
\node[cfone] (r) at (0,0) {Recommendation}; \node[cfdec] (q) at (3.8,0) {Eight questions};
\node[cfhope] (a) at (7.6,1.1) {Rely}; \node[cfthree] (e) at (7.6,0) {Escalate}; \node[cfrisk] (w) at (7.6,-1.1) {Withhold};
\draw[cfedge] (r)--(q); 
\draw[cfedge] (q)--(a.south west); 
\draw[cfedge] (q)--(e.west); 
\draw[cfedge] (q)--(w.north west);
\draw[cfedged] (e) to[out=180,in=-10] (q.south);
\end{cfwfig}
\vspace{-0.7cm}
\cfwcaption{Figure 11. The bedside trust decision: the eight questions
gate one decision with three guarded outcomes.}

\begin{cfwfig}
\draw[draw=black] (0,0) rectangle (6,4);
\draw[draw=black,dashed] (3,0)--(3,4); \draw[draw=black,dashed] (0,2)--(6,2);
\node[font=\scriptsize] at (3,-0.35) {demonstrated capability};
\node[font=\scriptsize,rotate=90] at (-0.35,2) {clinician reliance};
\node[cfhope,text width=14mm,minimum height=5mm] at (4.5,3.0) {Calibrated};
\node[cfrisk,text width=16mm,minimum height=5mm] at (1.5,3.0) {Automation bias};
\node[cfone,text width=14mm,minimum height=5mm] at (1.5,1.0) {Caution};
\node[cftwo,text width=16mm,minimum height=5mm] at (4.5,1.0) {Algorithm aversion};
\end{cfwfig}
\vspace{-0.7cm}
\cfwcaption{Figure 12. The calibrated-trust band: calibrated trust sits
between automation bias and algorithm aversion.}

\begin{cfwfig}
\node[cfone] (e) at (0,0) {Event}; \node[cftwo] (d) at (4.0,0) {Detect}; \node[cfdec] (g) at (7.5,0) {Gate 10};
\node[cfhope] (c) at (10.75,1.2) {Contain}; \node[cfrisk] (b) at (10.75,0) {BLOCK}; \node[cfthree] (s) at (10.75,-1.2) {Escalate};
\draw[cfedge] (e)--(d); \draw[cfedge] (d)--(g);
\draw[cfedge] (g)--(c.south west); \draw[cfedge] (g)--(b.west); \draw[cfedge] (g)--(s.north west);
\end{cfwfig}
\vspace{-0.7cm}
\cfwcaption{Figure 13. Safety containment: a detected event is gated
into containment, a hard block, or escalation.}

\begin{cfwfig}
\node[cfone] (a) at (0,0) {Review monitoring}; \node[cftwo] (b) at (4,0) {Receive advice};
\node[cfact] (c) at (8,0) {Accept or escalate}; \node[cfhope] (d) at (12,0) {Countersign record};
\draw[cfedge] (a)--(b); \draw[cfedge] (b)--(c); \draw[cfedge] (c)--(d);
\end{cfwfig}
\vspace{-0.7cm}
\cfwcaption{Figure 14. A clinician's day: the system enters the clinic
day at four scored touchpoints.}

\begin{cfwfig}
\node[cfone] (a) at (0,0) {Pilot}; \node[cftwo] (b) at (4,0) {Integrate};
\node[cfact] (c) at (8,0) {Operate}; \node[cfhope] (d) at (12,0) {Sustain};
\draw[cfedge] (a)--(b); \draw[cfedge] (b)--(c); \draw[cfedge] (c)--(d);
\end{cfwfig}
\vspace{-0.7cm}
\cfwcaption{Figure 15. The adoption maturity ladder: four phases from
pilot to standard work.}

\begin{cfwfig}
\node[cfthree] (y) at (0,0) {System}; \node[cfact] (o) at (3.8,0) {On-call clinician};
\node[cftwo] (a) at (7.6,0) {Attending}; \node[cfhope] (r) at (11.4,0) {Record};
\draw[cfedge] (y)--(o); \draw[cfedge] (o)--(a); \draw[cfedge] (a)--(r);
\draw[cfedge] (o) to[out=-60,in=-120,looseness=0.4] (r.south);
\end{cfwfig}
\vspace{-0.55cm}
\cfwcaption{Figure 16. The escalation pathway: an uncertain gate
escalates, and every step writes to the record.}

\vspace{0.1cm}

\begin{cfwfig}
\node[cftwo] (a) at (0,0) {Deployed}; \node[cfrisk] (b) at (4.0,0) {Drift detected};
\node[cfthree] (c) at (6.01,-1.2) {Retrain branch}; \node[cfact] (d) at (7.9,0) {Revalidated};
\node[cfhope] (e) at (11.7,0) {Redeployed};
\draw[cfedge] (a)--(b); \draw[cfedge] (b) to[out=-90,in=180] (c.west);
\draw[cfedge] (c) to[out=0,in=-90] (d.south); \draw[cfedge] (d)--(e);
\end{cfwfig}
\vspace{-0.55cm}
\cfwcaption{Figure 17. Model revalidation: a retrain branch revalidates
against the gates before it merges back into service.}

\vspace{0.1cm}

\begin{cfwfig}
\node[cfone] (p) at (0,0) {Patient 03:00}; \node[cftwo] (y) at (3.8,0) {System};
\node[cfdec] (g) at (7.0,0) {Ten gates}; \node[cfthree] (n) at (10.2,0) {Night clinician};
\draw[cfedge] (p)--(y); \draw[cfedge] (y)--(g); \draw[cfedge] (g)--(n);
\end{cfwfig}
\vspace{-0.55cm}
\cfwcaption{Figure 18. The night-shift scenario: 03:00 is treated
exactly like noon, by the same gates.}

\vspace{0.1cm}

\begin{cfwfig}
\node[cfact] (r) at (0,0) {Tumor board case}; \node[cfone] (c) at (4.0,1.2) {Validation};
\node[cftwo] (s) at (4.0,0.0) {Safety block}; \node[cfthree] (t) at (4.0,-1.2) {Record};
\node[cfhope] (e) at (8.0,0.0) {Subgroup and roles};
\draw[cfedge] (r.north east)--(c.south west); 
\draw[cfedge] (r.east)--(s.west); 
\draw[cfedge] (r.south east)--(t.north west); 
\draw[cfedge] (r.east) to[bend left=50, right=50, looseness=0.6] (e.west);
\end{cfwfig}
\vspace{-0.55cm}
\cfwcaption{Figure 19. Defensible at tumor board: the parallel supports
a clinician brings to peer review.}

\vspace{0.1cm}

\begin{cfwfig}
\node[cftwo] (r) at (0,0) {Running}; \node[cfact] (h) at (3.8,0) {Hold};
\node[cfrisk] (s) at (7.6,0) {Stopped};
\draw[cfedge] (r)--(h); \draw[cfedge] (h)--(s);
\draw[cfedged] (h) to[out=120,in=60, looseness=0.6] (r.north);
\draw[cfedge] (r) to[out=-40,in=-140, looseness=0.5] (s.south);
\end{cfwfig}
\vspace{-0.55cm}
\cfwcaption{Figure 20. Override and stop authority: a clinician can
pause, resume, or override to a full stop, and predicate can stop it.}

\vspace{0.2cm}
Taken together, the twenty figures trace one claim from a single bedside decision
(Figure 11) to trial-wide and platform trust (Figure 10), and each exhibit answers
one of the eight questions in the strict palette. None of them asks the clinician to
trust more than the evidence warrants; each shows where a gate, a record, a role, or
a control makes that trust calibrated.
